\documentclass[aps,pra,onecolumn,notitlepage,10pt,amsmath,amssymb,longbibliography,nofootinbib]{revtex4-2}
\usepackage[T1]{fontenc}
\usepackage{lmodern}
\usepackage{amsthm}
\usepackage{booktabs,longtable}
\usepackage[colorlinks=true,linkcolor=blue,citecolor=blue,urlcolor=blue]{hyperref}
\hypersetup{pdftitle={Maximum purity of absolutely PPT states of two qutrits},pdfauthor={Xiaojian Zeng}}
\renewcommand{\bibsection}{\section*{References}}
\newtheorem{theorem}{Theorem}
\newtheorem{lemma}{Lemma}
\newcommand{\Tr}{\operatorname{Tr}}
\newcommand{\APPT}{\operatorname{APPT}}
\newcommand{\R}{\mathbb{R}}
\newcommand{\C}{\mathbb{C}}
\newcommand{\lam}{\boldsymbol{\lambda}}
\begin{document}
\title{Maximum purity of absolutely PPT states of two qutrits}
\author{Xiaojian Zeng}
\email{xiaojianzeng@foxmail.com}
\affiliation{PBC School of Finance, Tsinghua University, 100190}
\date{September 12, 2026}
\begin{abstract}
We determine the maximum purity of a two-qutrit density operator that remains
positive under partial transpose after every global unitary transformation.
The exact value is $17/121$, and the unique decreasingly ordered maximizing
spectrum is $(3,1,1,1,1,1,1,1,1)/11$.
The computer-assisted proof uses three necessary linear inequalities from the
spectral PPT criterion. Their intersection with the ordered probability simplex
is an outer polytope with 49 vertices, all enumerated in exact rational arithmetic.
An active-constraint argument proves that the enumeration is complete, including
degenerate boundaries. The result gives the optimal purity-only sufficient
threshold for a global unitary to produce a negative partial transpose, and the
maximal Hilbert--Schmidt distance $4\sqrt{2}/33$ from the maximally mixed state
within the absolutely PPT set.
\end{abstract}
\maketitle

\section{Introduction}

A global unitary transformation preserves the spectrum and purity of a quantum
state, while it can change the state's entanglement across a fixed bipartition.
An absolutely positive-partial-transpose (APPT) state stays PPT along its entire
global unitary orbit. The largest purity of such a state therefore determines
when purity alone guarantees that some global unitary can produce a negative
partial transpose (NPT).

For two qutrits, Hildebrand's spectral criterion expresses APPT as two
$3\times3$ positive-semidefinite conditions~\cite{Hildebrand}.
Song and Chen gave an extreme-point criterion and identified the spectrum
$(3,1,1,1,1,1,1,1,1)/11$ among the two-eigenvalue extreme points; see Corollary~16,
Eq.~(10), of their preprint version~\cite{SongChen}.
Wang, Chen, and Song subsequently classified extreme points with exactly three
distinct eigenvalues~\cite{WangChenSong}. These restricted classifications do not
by themselves maximize purity over the full spectral domain.
Dung and Khoi reported the numerical estimate $0.1405$ for the two-qutrit APPT
maximum in Table~2 of their preprint~\cite{DungKhoi}.
Tran's Theorem~1 gives the upper bound $25/169$ in total dimension nine;
Remark~1.1 leaves the exact qutrit--qudit maximum open~\cite{Tran}.

The closest explicit conjecture is due to Ahiable, Kothakonda, and
Winter~\cite{AhiableKothakondaWinter}. Their Theorem~6.3 optimizes purity on an
\emph{inscribed} APPT polytope. For local dimensions $m=n=3$, its value is
$\max\{17/121,1/8\}=17/121$. Their Conjecture~6.7 asserts equality with the
maximum on the full APPT set. We prove precisely its $3\times3$ case.
The attaining spectrum is already known; the contribution here is the matching
upper bound on the entire domain and the uniqueness of the equality spectrum.
All preprint formula and theorem locators in this paper refer to the versions
specified in the bibliography.

\section{Model and sharp bound}
\label{sec:model}

Fix $\mathcal H=\C^3\otimes\C^3$ and let $\rho\succeq0$ with $\Tr\rho=1$.
Partial transpose $T_B$ is taken on the second factor in a fixed product basis.
The APPT condition is
\begin{equation}
 (U\rho U^\dagger)^{T_B}\succeq0\qquad\text{for every }U\in U(9).
 \label{eq:appt}
\end{equation}
There is no restriction on the global unitary, eigenvalue multiplicities, or
rank. Partial transpose is a mathematical test, not an assumed physical channel.
The state, its spectrum, and its purity are dimensionless. The problem is static
and finite dimensional; no perturbative expansion, dynamical approximation, or
spatial boundary condition is imposed.

Write the ordered spectrum as
\begin{equation}
 \lambda_1\geq\cdots\geq\lambda_9\geq0,
 \qquad \sum_{i=1}^9\lambda_i=1,
 \qquad P(\lam)=\Tr\rho^2=\sum_{i=1}^9\lambda_i^2.
 \label{eq:domain}
\end{equation}
By Eq.~(5) and Corollary~4 of the preprint of Ref.~\cite{Hildebrand},
Eq.~\eqref{eq:appt} is equivalent to both of the following matrices being
positive semidefinite:
\begin{align}
 L_1(\lam)&=\begin{pmatrix}
 2\lambda_9&\lambda_8-\lambda_1&\lambda_6-\lambda_2\\
 \lambda_8-\lambda_1&2\lambda_7&\lambda_5-\lambda_3\\
 \lambda_6-\lambda_2&\lambda_5-\lambda_3&2\lambda_4
 \end{pmatrix},\label{eq:L1}\\
 L_2(\lam)&=\begin{pmatrix}
 2\lambda_9&\lambda_8-\lambda_1&\lambda_7-\lambda_2\\
 \lambda_8-\lambda_1&2\lambda_6&\lambda_5-\lambda_3\\
 \lambda_7-\lambda_2&\lambda_5-\lambda_3&2\lambda_4
 \end{pmatrix}.\label{eq:L2}
\end{align}
Let $K$ denote the ordered spectral set satisfying
Eqs.~\eqref{eq:domain}--\eqref{eq:L2}, including both semidefinite conditions.

\begin{theorem}
\label{thm:main}
For every APPT state on $\C^3\otimes\C^3$,
\begin{equation}
 \Tr\rho^2\leq\frac{17}{121}.
 \label{eq:main}
\end{equation}
Equality holds if and only if the ordered spectrum is
\begin{equation}
 \lam_* = \frac1{11}(3,1,1,1,1,1,1,1,1).
 \label{eq:star}
\end{equation}
\end{theorem}

The proof below is computer assisted: a finite rational enumeration supplies
the vertex certificate, while an analytic argument establishes its completeness.
The complete certificate is printed in Appendix~\ref{app:vertices}.

\section{Three necessary inequalities}
\label{sec:cuts}

For $v_1=(1,1,0)^T$, $v_2=(1,1,1)^T$, and $v_3=(2,1,0)^T$, define
$g_j(\lam)=v_j^TL_1(\lam)v_j/2$. Direct expansion gives
\begin{align}
 g_1(\lam)&=-\lambda_1+\lambda_7+\lambda_8+\lambda_9,\label{eq:g1}\\
 g_2(\lam)&=-\lambda_1-\lambda_2-\lambda_3+
              \lambda_4+\lambda_5+\lambda_6+\lambda_7+\lambda_8+\lambda_9,
              \label{eq:g2}\\
 g_3(\lam)&=-2\lambda_1+\lambda_7+2\lambda_8+4\lambda_9.\label{eq:g3}
\end{align}
Every $\lam\in K$ obeys $g_j(\lam)\geq0$. In particular,
$g_2\geq0$ is equivalent to $\lambda_1+\lambda_2+\lambda_3\leq1/2$.

Define the outer polytope
\begin{equation}
 Q=\{\lam\text{ satisfying Eq.~\eqref{eq:domain}}:
                    g_1(\lam),g_2(\lam),g_3(\lam)\geq0\}.
 \label{eq:Q}
\end{equation}
Then $K\subseteq Q$. Only this inclusion is needed for the upper bound; the
three inequalities are not asserted to be a sufficient APPT criterion.
In particular, $L_2\succeq0$ remains part of the physical domain and will also
be checked at the equality spectrum.

The same inequalities can be checked directly from Schmidt coefficients.
For $b=\sum_{i=1}^3 x_i|ii\rangle$, the eigenvalues of
$(|b\rangle\langle b|)^{T_B}$ are $x_i^2$ and $\pm x_ix_j$ for $i<j$.
Pairing this spectrum in increasing order with $\lam$ in decreasing order
gives the minimum trace over its unitary orbit~\cite{Hildebrand}.
The choices $x=(1,1,0),(1,1,1),(2,1,0)$ respectively give the ordered spectra
\begin{equation}
 (-1,0,0,0,0,0,1,1,1),\quad
 (-1,-1,-1,1,1,1,1,1,1),\quad
 (-2,0,0,0,0,0,1,2,4),
 \label{eq:schmidtchecks}
\end{equation}
and hence the three linear forms above. Normalization of $b$ is unnecessary
because positive scaling does not change the sign of the test.

\section{Complete exact vertex certificate}
\label{sec:certificate}

\subsection{Completeness of the enumeration}

Set $\lambda_{10}=0$ and introduce gap weights
\begin{equation}
 y_k=k(\lambda_k-\lambda_{k+1}),\qquad
 \lambda_i=\sum_{k=i}^9\frac{y_k}{k}.
 \label{eq:gaps}
\end{equation}
This is a linear bijection between the full ordered probability simplex and
$\{y\in\R^9:y\geq0,\ \boldsymbol{1}^Ty=1\}$.
It is the standard representation by the uniform rank-$k$ spectra, also used
in Ref.~\cite{AhiableKothakondaWinter}; it imposes no multiplicity assumption.
In these coordinates the three extra inequalities are $Ay\geq0$, where
\begin{equation}
 A=\begin{pmatrix}
 -1&-\frac12&-\frac13&-\frac14&-\frac15&-\frac16&0&\frac18&\frac29\\[2pt]
 -1&-1&-1&-\frac12&-\frac15&0&\frac17&\frac14&\frac13\\[2pt]
 -2&-1&-\frac23&-\frac12&-\frac25&-\frac13&-\frac17&\frac18&\frac59
 \end{pmatrix}.
 \label{eq:A}
\end{equation}

\begin{lemma}
\label{lem:complete}
Each vertex of $\{y\geq0:\boldsymbol{1}^Ty=1,Ay\geq0\}$ has at most four
positive coordinates. If its positive support is $S$, with $s=|S|$, some
$s-1$ active rows of $A$, together with the trace row, form a nonsingular
$s\times s$ system on $S$ whose solution is the vertex.
\end{lemma}
\begin{proof}
Within the simplex face with positive support $S$, the space of directions
$d$ supported on $S$ and satisfying $\boldsymbol{1}^Td=0$ has dimension $s-1$.
If the active rows of $A$ do not eliminate this space, there is a nonzero
direction $d$ that also annihilates all active rows. For sufficiently small
$\epsilon>0$, both $y+\epsilon d$ and $y-\epsilon d$ preserve the positive
coordinates and all strictly slack inequalities. They are distinct feasible
points with midpoint $y$, contradicting vertexhood. Thus $s-1$ independent
active restrictions exist. Since $A$ has only three rows, $s\leq4$; choosing
such restrictions gives the stated nonsingular system.
\end{proof}

Consequently, every vertex is found by considering $s=1,\ldots,4$, every
support $S\subseteq\{1,\ldots,9\}$ of size $s$, and every subset $J$ of
$s-1$ rows of $A$. For each nonsingular matrix
\begin{equation}
 M_{S,J}=\begin{pmatrix}\boldsymbol{1}_S^T\\A_{J,S}\end{pmatrix},
 \qquad M_{S,J}y_S=(1,0,\ldots,0)^T,
 \label{eq:basis}
\end{equation}
solve exactly, set $y_{S^c}=0$, and check all inequalities $y\geq0$ and
$Ay\geq0$. Each feasible solution is a vertex: the selected independent
equalities and the zero coordinates determine it uniquely. Solutions with zero
weights are retained and deduplicated, so degeneracies are included.
There are exactly
\begin{equation}
 \sum_{s=1}^4 \binom{9}{s}\binom{3}{s-1}
       =9+108+252+126=495
 \label{eq:count}
\end{equation}
systems to consider. Singular systems may be discarded because
Lemma~\ref{lem:complete} supplies a nonsingular basis for every vertex.
The bound on positive gaps concerns vertices of the outer polytope and is a
consequence of its three cuts, not a restriction on arbitrary APPT states.

\subsection{Rational evaluation and cross-checks}

Exact enumeration gives 72 singular systems and 423 nonsingular systems.
There are 63 feasible bases, which yield 49 distinct vertices after removing
duplicates. Table~\ref{tab:groups} summarizes the results by the actual positive
support size. Every vertex and its exact purity are listed in
Appendix~\ref{app:vertices}. Purity can be evaluated directly from the listed
weights using
\begin{equation}
 P(y)=\sum_{k,l=1}^9\frac{y_ky_l}{\max(k,l)},
 \label{eq:Pgap}
\end{equation}
which follows by substituting Eq.~\eqref{eq:gaps} into $\sum_i\lambda_i^2$.

\begin{table}[htb]
\caption{Complete vertex enumeration for $Q$, grouped by the number $s$ of
positive gap weights. These are exact rational results.}
\label{tab:groups}
\begin{ruledtabular}
\begin{tabular}{ccc}
$s$ & Number of vertices & Largest purity in the group\\
1&2&$1/8$\\
2&14&$17/121$\\
3&21&$5/36$\\
4&12&$3/22$
\end{tabular}
\end{ruledtabular}
\end{table}

The unique vertex with purity $17/121$ has
$y_1=2/11$, $y_9=9/11$, and all other gap weights zero.
Every other vertex has purity at most $5/36$, with a strict gap
\begin{equation}
 \frac{17}{121}-\frac5{36}=\frac7{4356}>0.
 \label{eq:gap}
\end{equation}

For a separate representation of the finite calculation, one may work directly
in the nine spectrum coordinates. There are twelve homogeneous inequalities:
eight ordering inequalities, $\lambda_9\geq0$, and the three $g_j\geq0$.
Enumerating all $\binom{12}{8}=495$ choices of eight active inequalities,
together with the trace row, and solving by rational Gauss--Jordan elimination
reproduces all 49 vertices. The supplied standard-library Python program
implements this calculation and checks the complete distributed table only
after generating its vertex set. A separate exact incremental clipping
calculation also reproduces the same set; adding $g_1,g_2,g_3$ in that order
gives 15, 31, and 49 vertices and maxima $25/169$, $1/7$, and $17/121$.
The first step recovers the dimension-nine relaxation underlying
Ref.~\cite{Tran}.

All decisions in the certificate use rational arithmetic, with no feasibility
tolerance. Floating-point optimization is not part of the proof. Agreement
between implementations is a cross-check; completeness rests on the
active-constraint argument above. Regeneration instructions and the exact
data accompany the manuscript.

\subsection{Research tools and verification}

Generative AI (OpenAI Codex) was used substantively for literature exploration,
selection of test vectors, development of the proof, preparation and execution
of exact-arithmetic code, internal checks, and manuscript drafting.
The computational checks consisted of complete rational enumerations in
different representations and a separate rational clipping implementation.
These checks do not constitute human peer review. This local draft still
requires the named author's scientific review and final approval; no completed
human verification is asserted.

\section{Proof of the sharp bound and its consequences}

\begin{proof}[Proof of Theorem~\ref{thm:main}]
The compact polytope $Q$ is the convex hull of its vertices. Since
$P(\lam)=\|\lam\|_2^2$ is convex, the exact vertex bound implies
$P\leq17/121$ on all of $Q$, hence on $K\subseteq Q$.
Strict convexity rules out a nonvertex maximizer: a nontrivial convex
combination of distinct feasible points has purity strictly below the
corresponding convex combination of their purities. The only possible ordered
equality spectrum is therefore the unique maximizing vertex $\lam_*$.

At this spectrum, both complete semidefinite constraints become
\begin{equation}
 L_1(\lam_*)=L_2(\lam_*)=
 \frac2{11}\begin{pmatrix}1&-1&0\\-1&1&0\\0&0&1\end{pmatrix}\succeq0.
 \label{eq:attain}
\end{equation}
Their eigenvalues are $0,2/11,4/11$. Thus $\lam_*\in K$, and
$P(\lam_*)=(9+8)/121=17/121$, proving attainment and uniqueness.
\end{proof}

Attainment also has a direct full-orbit check. Every state with spectrum
$\lam_*$ has the form
\begin{equation}
 \rho=\frac{I_9+2|\psi\rangle\langle\psi|}{11}
 \label{eq:orbit}
\end{equation}
for a normalized $|\psi\rangle$. If its Schmidt coefficients are $x_i$, the
potentially smallest eigenvalues of $\rho^{T_B}$ are
$(1-2x_ix_j)/11$. They are nonnegative because
$2x_ix_j\leq x_i^2+x_j^2\leq1$.

For any two-qutrit state, Theorem~\ref{thm:main} implies
\begin{equation}
 \Tr\rho^2>\frac{17}{121}\quad\Longrightarrow\quad
 \exists U\in U(9): (U\rho U^\dagger)^{T_B}\not\succeq0.
 \label{eq:NPT}
\end{equation}
This purity-only sufficient threshold cannot be lowered, since the APPT orbit
in Eq.~\eqref{eq:orbit} reaches it. The implication is not a characterization
of states at lower purity: the theorem does not assert that every state below
the threshold is APPT. It also does not assert distillability of the resulting
NPT state or identify APPT with absolute separability.

Finally, for the Hilbert--Schmidt norm
$\|X\|_{\mathrm{HS}}=(\Tr X^\dagger X)^{1/2}$,
\begin{equation}
 \max_{\rho\in\APPT_{3,3}}
    \left\|\rho-\frac{I_9}{9}\right\|_{\mathrm{HS}}^2
       =\frac{17}{121}-\frac19=\frac{32}{1089}.
 \label{eq:radius}
\end{equation}
The maximum distance is therefore $4\sqrt2/33$, attained precisely by the
unitary orbit in Eq.~\eqref{eq:orbit}.

\section{Discussion}

Three necessary linear tests suffice to enclose the full two-qutrit APPT
spectral domain in a polytope with the correct largest purity. The exact
upper bound improves $25/169$ by $152/20449$ and matches the previously known
attaining spectrum. This establishes the $3\times3$ instance of the
maximum-purity conjecture of Ref.~\cite{AhiableKothakondaWinter}, including
its maximizing spectrum. No conclusion for general $3\times n$ or other
local dimensions follows from this nine-variable certificate.

The argument avoids a classification of all APPT extreme points: it classifies
only the vertices of a deliberately larger set. Its limitation is equally
specific. The finite rational table and the analytic completeness argument
together give a computer-assisted proof, rather than a floating-point search
or a proof-assistant formalization. An analytic compression of the certificate
could simplify the presentation, but is not assumed here.

\section*{Data availability}
All exact vertex data and a Python standard-library regeneration program are
included with this local manuscript in the accompanying reproducibility
materials. The full vertex table is also printed below. No experimental or
third-party numerical dataset is used, and no public repository deposit is
claimed.

\appendix
\section{Complete rational vertex certificate}
\label{app:vertices}

Table~\ref{tab:vertices} gives all vertices of $Q$. An entry $k:a/b$ means
$y_k=a/b$; every unlisted weight is zero. The ordered spectrum is recovered
from Eq.~\eqref{eq:gaps}. Thus the table, Eq.~\eqref{eq:A}, and
Eq.~\eqref{eq:Pgap} give a finite exact certificate whose entries can be checked
without numerical tolerances. Table numbering fixes one ordering of tied
purities and has no mathematical significance.

\begingroup
\small
\setlength{\tabcolsep}{10pt}
\begin{longtable}{rll}
\caption{All 49 vertices of the three-cut outer polytope, in gap weights.}
\label{tab:vertices}\\
\toprule
No. & Positive gap weights $k:y_k$ & $P$\\
\midrule
\endfirsthead
\multicolumn{3}{c}{Table \thetable\ continued}\\
\toprule
No. & Positive gap weights $k:y_k$ & $P$\\
\midrule
\endhead
\midrule
\multicolumn{3}{r}{Continued on next page}\\
\endfoot
\bottomrule
\endlastfoot
1 & $1:2/11,\; 9:9/11$ & $17/121$\\
2 & $1:1/12,\; 2:1/6,\; 9:3/4$ & $5/36$\\
3 & $1:1/22,\; 2:2/11,\; 8:4/11,\; 9:9/22$ & $3/22$\\
4 & $2:7/32,\; 8:1/2,\; 9:9/32$ & $279/2048$\\
5 & $4:2/9,\; 5:5/18,\; 9:1/2$ & $11/81$\\
6 & $5:10/19,\; 9:9/19$ & $49/361$\\
7 & $4:4/15,\; 5:1/6,\; 8:4/15,\; 9:3/10$ & $61/450$\\
8 & $4:14/39,\; 8:16/39,\; 9:3/13$ & $206/1521$\\
9 & $4:1/3,\; 7:7/24,\; 9:3/8$ & $13/96$\\
10 & $2:1/4,\; 9:3/4$ & $13/96$\\
11 & $1:1/8,\; 3:1/8,\; 9:3/4$ & $13/96$\\
12 & $4:1/3,\; 6:1/10,\; 8:4/15,\; 9:3/10$ & $27/200$\\
13 & $2:1/5,\; 7:7/20,\; 9:9/20$ & $27/200$\\
14 & $1:3/20,\; 8:2/5,\; 9:9/20$ & $27/200$\\
15 & $5:15/32,\; 8:1/4,\; 9:9/32$ & $69/512$\\
16 & $2:1/5,\; 6:3/25,\; 8:8/25,\; 9:9/25$ & $84/625$\\
17 & $3:1/12,\; 5:5/12,\; 9:1/2$ & $29/216$\\
18 & $4:1/3,\; 6:1/6,\; 9:1/2$ & $29/216$\\
19 & $3:1/10,\; 5:1/3,\; 8:4/15,\; 9:3/10$ & $2/15$\\
20 & $4:2/5,\; 9:3/5$ & $2/15$\\
21 & $2:1/5,\; 6:1/5,\; 9:3/5$ & $2/15$\\
22 & $5:5/13,\; 7:7/26,\; 9:9/26$ & $45/338$\\
23 & $1:1/26,\; 4:4/13,\; 8:4/13,\; 9:9/26$ & $45/338$\\
24 & $2:2/13,\; 5:5/26,\; 8:4/13,\; 9:9/26$ & $45/338$\\
25 & $2:1/8,\; 5:5/16,\; 9:9/16$ & $17/128$\\
26 & $1:1/14,\; 4:2/7,\; 9:9/14$ & $13/98$\\
27 & $1:1/11,\; 3:3/22,\; 8:4/11,\; 9:9/22$ & $16/121$\\
28 & $3:1/8,\; 5:5/24,\; 7:7/24,\; 9:3/8$ & $19/144$\\
29 & $3:7/34,\; 8:12/17,\; 9:3/34$ & $38/289$\\
30 & $7:35/44,\; 9:9/44$ & $127/968$\\
31 & $3:2/11,\; 7:21/44,\; 9:15/44$ & $127/968$\\
32 & $3:1/6,\; 6:1/6,\; 7:7/24,\; 9:3/8$ & $113/864$\\
33 & $3:1/6,\; 6:4/15,\; 8:4/15,\; 9:3/10$ & $88/675$\\
34 & $3:3/19,\; 8:16/19$ & $47/361$\\
35 & $6:3/7,\; 7:1/4,\; 9:9/28$ & $51/392$\\
36 & $4:1/5,\; 8:4/5$ & $13/100$\\
37 & $1:1/20,\; 3:3/20,\; 7:7/20,\; 9:9/20$ & $13/100$\\
38 & $6:18/35,\; 8:8/35,\; 9:9/35$ & $159/1225$\\
39 & $2:1/9,\; 8:8/9$ & $7/54$\\
40 & $3:1/6,\; 6:1/3,\; 9:1/2$ & $7/54$\\
41 & $1:1/9,\; 7:7/18,\; 9:1/2$ & $7/54$\\
42 & $5:5/21,\; 8:16/21$ & $19/147$\\
43 & $6:4/7,\; 9:3/7$ & $19/147$\\
44 & $7:7/15,\; 8:8/15$ & $29/225$\\
45 & $6:3/11,\; 8:8/11$ & $31/242$\\
46 & $1:1/17,\; 8:16/17$ & $37/289$\\
47 & $8:1$ & $1/8$\\
48 & $3:1/4,\; 9:3/4$ & $1/8$\\
49 & $9:1$ & $1/9$\\
\end{longtable}
\endgroup


% Load the formatting control; all six literature entries are cited in the text.
\nocite{*}
\bibliographystyle{apsrev4-2}
\raggedright
\bibliography{references}
\end{document}
